\chapter{Localizzazione continua di un servizio}
\label{cap:localizzazione}

\noindent\textbf{Classe:} NLP convesso \hfill
\textbf{Script:} \texttt{python/lab09\_localizzazione.py}

\section{Motivazione gestionale}

Dove collocare una stazione di ricarica rapida, un micro-hub logistico, un presidio
sanitario, per essere ``vicini'' alla domanda? La risposta dipende da cosa significa
vicini: distanza \emph{media} (efficienza), distanza \emph{massima} (equità verso
l'utente più sfortunato) o distanza \emph{quadratica} (che porta al baricentro).
Tre obiettivi, tre punti diversi sulla mappa: la geometria rende il trade-off visibile
come in nessun altro modello.



Le decisioni di localizzazione sono tra le più costose da correggere: un hub nel
posto sbagliato resta nel posto sbagliato per anni. La versione continua del
problema (coordinate libere sul piano) è il laboratorio ideale per un tema che
riappare in tutte le decisioni pubbliche e private: il conflitto tra
\emph{efficienza} (servire bene in media) ed \emph{equità} (non abbandonare
nessuno).

\textbf{In questo capitolo impareremo a:} (1) formulare tre obiettivi di
localizzazione e capire quando coincidono; (2) risolvere NLP con scipy da buoni
punti iniziali; (3) gestire vincoli geografici convessi; (4) costruire una
frontiera efficienza-equità con il metodo $\varepsilon$-constraint.

\section{Costruzione guidata del modello}

\paragraph{Il problema a parole.} \emph{Decidiamo} le coordinate in cui collocare la
nuova struttura. \emph{L'obiettivo} --- ed è la vera scelta manageriale --- può essere
minimizzare la distanza media pesata dagli utenti, oppure la distanza dell'utente più
lontano, oppure la distanza quadratica. \emph{I vincoli}, quando presenti, delimitano
l'area ammissibile (una zona rettangolare, un raggio massimo da un'infrastruttura).

\subsection*{Insiemi e dati (input del modello)}
\begin{center}\small
\begin{tabular}{llp{7.2cm}}
\toprule
Simbolo & Tipo & Significato \\
\midrule
$I$ & insieme finito & punti di domanda (quartieri) \\
$(a_i, b_i)$ & $\in \R^2$ & coordinate del quartiere $i \in I$ \\
$w_i$ & $\in \R_{\ge 0}$ & peso del quartiere: popolazione, domanda o priorità \\
\bottomrule
\end{tabular}
\end{center}

\subsection*{Variabili decisionali}
\begin{center}\small
\begin{tabular}{llp{7.2cm}}
\toprule
$(x, y)$ & $\in \R^2$, continue & coordinate della nuova struttura \\
\bottomrule
\end{tabular}
\end{center}

\begin{modello}[Tre obiettivi classici]
\begin{align*}
\text{Weber:}\quad & \min_{x,y}\; \textstyle\sum_{i \in I} w_i \sqrt{(x-a_i)^2 + (y-b_i)^2}\\
\text{Minimax:}\quad & \min z \;\;\text{s.t.}\;\; \sqrt{(x-a_i)^2 + (y-b_i)^2} \le z \;\;\forall i \in I\\
\text{Quadratico:}\quad & \min_{x,y}\; \textstyle\sum_{i \in I} w_i \bigl[(x-a_i)^2 + (y-b_i)^2\bigr]
\end{align*}
Tutti convessi (norme e quadrati di funzioni affini). Vincoli geografici convessi
tipici: zona rettangolare $x^L \le x \le x^U$; distanza massima da un'infrastruttura
$(x - a_0)^2 + (y - b_0)^2 \le r^2$, con $(a_0, b_0) \in \R^2$ e raggio $r > 0$ dati.
\end{modello}

\begin{spiegazione}
\textbf{Quadratico $\Rightarrow$ baricentro.} Annullando il gradiente:
$x^* = \sum_{i \in I} w_i a_i \big/ \sum_{i \in I} w_i$ (e analogo per $y$): l'unico caso con formula
chiusa, utile come punto di partenza per gli altri.

\textbf{Weber.} Non ha forma chiusa; la funzione non è differenziabile nei punti
$(a_i, b_i)$. All'ottimo i ``tiri'' unitari pesati dei quartieri si annullano:
$\sum_{i \in I} w_i\, \vet u_i = \vet 0$ dove $\vet u_i$ è il versore dalla struttura verso il quartiere $i$.

\textbf{Minimax.} L'ottimo è il centro del \emph{cerchio minimo} che racchiude tutti i
punti: dipende solo dai quartieri estremi e ignora i pesi --- per questo protegge le
periferie.
\end{spiegazione}

\section{Esempio svolto a mano}

\begin{esempio}[Tre clienti su una strada]
Clienti nei km $0$, $4$ e $10$ con pesi $1$, $1$ e $3$.

\textbf{Weber 1D $=$ mediana pesata.} Il costo è
$f(x) = |x| + |x - 4| + 3|x - 10|$. Spostandosi di $\delta$ a sinistra da $x = 10$,
il costo dei primi due clienti scende di $2\delta$ ma quello del terzo sale di
$3\delta$: peggiora. Da destra, simmetricamente, peggiora ancora. Ottimo: $x^* = 10$,
il punto che lascia \emph{almeno metà del peso totale} ($5/2 = 2{,}5$) da entrambi i
lati --- la mediana pesata.

\textbf{Baricentro.} $x = (1 \cdot 0 + 1 \cdot 4 + 3 \cdot 10)/5 = 6{,}8$:
sensibile a \emph{quanto} sono lontani i punti, non solo da che parte stanno.

\textbf{Minimax.} Punto medio degli estremi: $(0 + 10)/2 = 5$, pesi ignorati.

Tre obiettivi ragionevoli, tre risposte diverse: $10$, $6{,}8$, $5$. La scelta
dell'obiettivo È la decisione manageriale.
\end{esempio}

\section{Caso di studio}

Dodici quartieri con pesi tra 5 e 25 mila abitanti
(\texttt{dati/localizzazione\_quartieri.csv}); vincolo tecnico: la stazione deve
stare entro $r = 2$ km dalla cabina elettrica in $(7, 6)$.

\begin{lstlisting}[style=python]
res = minimize(f_weber, x0=baricentro, method="Nelder-Mead")
res_m = minimize(f_max, x0=baricentro, method="Nelder-Mead")
# con vincolo geografico (SLSQP, vincolo >= 0)
res_v = minimize(f_weber, x0=cabina, method="SLSQP",
    constraints=[{"type": "ineq",
                  "fun": lambda p: R**2 - ((p - cabina)**2).sum()}])
\end{lstlisting}

\section{Risultati}

\begin{lstlisting}[style=output]
Baricentro pesato : (4,897, 5,397)
Weber             : (4,886, 5,310)  dist. media 3,344 km  max 6,314 km
Minimax           : (5,496, 5,029)  dist. media 3,391 km  max 5,680 km
Weber vincolato   : (5,083, 5,429)  costo +0,2% (vincolo attivo, ottimo sul bordo)
\end{lstlisting}

\begin{center}
% Mappa dei quartieri e localizzazioni ottime (generato da lab09_localizzazione.py)
\begin{tikzpicture}[scale=0.82, >=stealth]
  \draw[black!20, very thin] (0,0) grid[step=1] (10.5,10);
  \draw[->, black!50] (0,0) -- (10.8,0) node[below left, font=\scriptsize] {km est};
  \draw[->, black!50] (0,0) -- (0,10.3) node[below left, rotate=90, font=\scriptsize] {km nord};
  \fill[teal, opacity=0.45] (1.00,8.50) circle (0.38);
  \node[font=\tiny, text=black!55, anchor=west] at (1.38,8.50) {Q1};
  \fill[teal, opacity=0.45] (2.50,6.00) circle (0.47);
  \node[font=\tiny, text=black!55, anchor=west] at (2.97,6.00) {Q2};
  \fill[teal, opacity=0.45] (4.00,9.00) circle (0.33);
  \node[font=\tiny, text=black!55, anchor=west] at (4.33,9.00) {Q3};
  \fill[teal, opacity=0.45] (5.50,7.50) circle (0.52);
  \node[font=\tiny, text=black!55, anchor=west] at (6.02,7.50) {Q4};
  \fill[teal, opacity=0.45] (7.00,8.00) circle (0.43);
  \node[font=\tiny, text=black!55, anchor=west] at (7.43,8.00) {Q5};
  \fill[teal, opacity=0.45] (9.00,9.50) circle (0.27);
  \node[font=\tiny, text=black!55, anchor=west] at (9.27,9.50) {Q6};
  \fill[teal, opacity=0.45] (1.50,3.00) circle (0.41);
  \node[font=\tiny, text=black!55, anchor=west] at (1.91,3.00) {Q7};
  \fill[teal, opacity=0.45] (3.00,1.50) circle (0.31);
  \node[font=\tiny, text=black!55, anchor=west] at (3.31,1.50) {Q8};
  \fill[teal, opacity=0.45] (5.00,3.50) circle (0.55);
  \node[font=\tiny, text=black!55, anchor=west] at (5.55,3.50) {Q9};
  \fill[teal, opacity=0.45] (6.50,2.00) circle (0.35);
  \node[font=\tiny, text=black!55, anchor=west] at (6.85,2.00) {Q10};
  \fill[teal, opacity=0.45] (8.00,4.00) circle (0.44);
  \node[font=\tiny, text=black!55, anchor=west] at (8.44,4.00) {Q11};
  \fill[teal, opacity=0.45] (9.50,1.00) circle (0.25);
  \node[font=\tiny, text=black!55, anchor=west] at (9.75,1.00) {Q12};
  % traiettoria del compromesso efficienza-equità
  \draw[black!45, thick, densely dotted] (5.496,5.029) -- (5.449,5.026) -- (5.401,5.024) -- (5.354,5.021) -- (5.308,5.018) -- (5.261,5.015) -- (5.215,5.013) -- (5.169,5.010) -- (5.124,5.007) -- (5.078,5.005) -- (5.048,5.018) -- (5.032,5.048) -- (5.016,5.077) -- (5.000,5.107) -- (4.983,5.136) -- (4.967,5.165) -- (4.951,5.194) -- (4.935,5.223) -- (4.918,5.253) -- (4.902,5.282) -- (4.886,5.310);
  % vincolo geografico: cerchio della cabina
  \draw[viola, dashed, thick] (7.0,6.0) circle (2.0);
  \node[rectangle, fill=viola, minimum size=2.2mm, inner sep=0] at (7.0,6.0) {};
  \node[star, star points=5, fill=rossomattone, minimum size=3.2mm, inner sep=0pt, label={[font=\scriptsize, text=rossomattone, label distance=2.5mm]190:Weber}] at (4.886,5.310) {};
  \node[star, star points=5, fill=arancio, minimum size=3.2mm, inner sep=0pt, label={[font=\scriptsize, text=arancio, label distance=2.5mm]-55:minimax}] at (5.496,5.029) {};
  \node[star, star points=5, fill=verde, minimum size=3.2mm, inner sep=0pt, label={[font=\scriptsize, text=verde, label distance=2.5mm]120:baricentro}] at (4.897,5.397) {};
  \node[star, star points=5, fill=viola, minimum size=3.2mm, inner sep=0pt, label={[font=\scriptsize, text=viola, label distance=2.5mm]35:Weber vincolato}] at (5.083,5.429) {};
\end{tikzpicture}
\captionof{figure}{Mappa della città: i cerchi teal sono i 12 quartieri (area
proporzionale alla popolazione); le stelle indicano le quattro localizzazioni ottime
(Weber in rosso, minimax in arancione, baricentro in verde, Weber vincolato in
viola); il cerchio viola tratteggiato è il vincolo dei 2 km dalla cabina elettrica;
la linea punteggiata è la traiettoria delle soluzioni di compromesso tra minimax e
Weber.}
\end{center}

\begin{insight}
Weber e baricentro quasi coincidono qui (i pesi sono distribuiti in modo bilanciato),
ma il minimax si sposta di oltre mezzo km verso sud-est per proteggere i quartieri
periferici. Il vincolo della cabina costa pochissimo (+0,2\% di distanza totale):
scoprire che un vincolo temuto è quasi gratis è un risultato manageriale importante
quanto l'ottimo stesso.
\end{insight}

\section{Analisi di sensitività: la frontiera efficienza-equità}

Minimizziamo la distanza media imponendo un tetto $D$ alla distanza massima
(metodo $\varepsilon$-constraint), e variamo $D$ tra i due estremi:

\begin{center}
\begin{tikzpicture}
\begin{axis}[lab, title={Frontiera efficienza-equità},
  xlabel={distanza massima consentita $D$ (km) --- equità},
  ylabel={distanza media pesata (km)}]
\addplot+ [mark=*, mark size=1.5pt] table [col sep=comma, x=dist_max, y=dist_media]
  {\figdat{cap09_frontiera}};
\node [font=\scriptsize, anchor=south west] at (axis cs:5.70,3.388) {minimax};
\node [font=\scriptsize, anchor=north east] at (axis cs:6.31,3.346) {Weber};
\end{axis}
\end{tikzpicture}
\captionof{figure}{Frontiera efficienza-equità ($\varepsilon$-constraint): distanza
media pesata minima ottenibile per ogni tetto $D$ imposto alla distanza massima.
Gli estremi sono la soluzione minimax (a sinistra) e il punto di Weber (a destra);
la curva quasi piatta indica che l'equità costa pochissimo in efficienza.}
\end{center}

\begin{lstlisting}[style=output]
tetto D = 5,680 km: media 3,391   (soluzione minimax)
tetto D = 5,997 km: media 3,351
tetto D = 6,314 km: media 3,344   (soluzione Weber)
\end{lstlisting}

\begin{spiegazione}
La frontiera è quasi piatta: garantire al quartiere più lontano $0{,}63$ km in meno
costa solo $47$ metri di distanza media. Quando la frontiera è così piatta, la
soluzione equa è ``quasi gratis'' --- un argomento fortissimo in una discussione
pubblica. Il caso opposto (frontiera ripida) segnalerebbe un vero conflitto
efficienza-equità.
\end{spiegazione}

\section{Esercizi}

\begin{esercizio}[verifica]
Calcolare a mano il baricentro pesato dei 12 quartieri dai dati CSV e confrontarlo
con l'output dello script.
\end{esercizio}

\begin{esercizio}[peso che cambia tutto]
Portare il peso di Q12 (il quartiere in basso a destra) da 5 a 40: di quanto si
spostano Weber e minimax? Quale dei due è più sensibile ai pesi, e perché?
\end{esercizio}

\begin{esercizio}[due strutture]
Con \emph{due} stazioni e l'assegnazione di ogni quartiere alla più vicina, il
problema resta convesso? Provare un approccio alternato (assegna--ottimizza) e
discutere perché trova solo ottimi locali.
\end{esercizio}

\begin{esercizio}[distanza Manhattan]
Sostituire la distanza euclidea con $|x - a_i| + |y - b_i|$ e mostrare che il
problema di Weber diventa un LP (introdurre variabili ausiliarie per i valori
assoluti). Risolverlo con Gurobi e confrontare l'ottimo.
\end{esercizio}

\begin{esercizio}[moltiplicatore del raggio]
Stimare per perturbazione il moltiplicatore del vincolo $r = 2$ km (risolvere con
$r = 2{,}1$): quanto vale, in km di distanza totale, ogni 100 m di raggio in più?
\end{esercizio}
